% =============================================================================
% BIBLIOGRAPHIE (texte APA)
% Source : consigne/revue_litterature_v2_DSR.tex (lignes 2157 a 2563)
% Stockee comme bibliographie manuelle en attendant la conversion BibTeX
% progressive vers bibliography.bib + biblatex-apa.
% =============================================================================

\chapter*{Bibliographie}
\addcontentsline{toc}{chapter}{Bibliographie}

\begingroup
\setlength{\parskip}{0.6em}
\setlength{\parindent}{-1.5em}
\leftskip=1.5em\noindent



Amazon. (2024). \emph{Amazon Q Developer Agent: Java migration
capabilities}. \url{https://aws.amazon.com/q/developer/}

Amazon. (2025). \emph{PolyMigration: A benchmark for cross-language code
migration}. Amazon Web Services.
\url{https://aws.amazon.com/q/developer/polymigration/}

Anthropic. (2025). \emph{When to use multi-agent systems (and when not
to)}. Claude Documentation. \url{https://docs.anthropic.com/}

Anthropic. (2025b, December 2). \emph{How AI is transforming work at
Anthropic}.
\url{https://www.anthropic.com/news/how-ai-is-transforming-work-at-anthropic}

Anthropic. (2025c). \emph{Model Context Protocol (MCP)}.
\url{https://modelcontextprotocol.io/}

Anthropic. (2026). \emph{Claude Code}. \url{https://claude.ai/code}

Barke, S., James, M. B., \& Polikarpova, N. (2023). Grounded Copilot:
How programmers interact with code-generating models. \emph{Proceedings
of the ACM on Programming Languages, 7}(OOPSLA1), 85-111.
\url{https://doi.org/10.1145/3586030}

Berti, A., Maatallah, M., Jessen, U., Sroka, M., \& Ghannouchi, S. A.
(2024). Re-thinking process mining in the AI-based agents era.
\emph{arXiv preprint arXiv:2408.07720}.
\url{https://doi.org/10.48550/arXiv.2408.07720}

Bolici, F., Howison, J., \& Crowston, K. (2016). Stigmergic coordination
in FLOSS development teams: Integrating explicit and implicit
mechanisms. \emph{Cognitive Systems Research, 38}, 14-22.
\url{https://doi.org/10.1016/j.cogsys.2015.12.003}

Bonabeau, E., Dorigo, M., \& Theraulaz, G. (1999). \emph{Swarm
intelligence: From natural to artificial systems}. Oxford University
Press.

Carriero, N., \& Gelernter, D. (1992). Coordination languages and their
significance. \emph{Communications of the ACM, 35}(2), 97-107.
\url{https://doi.org/10.1145/129630.129635}

Cemri, M., Pan, M. Z., Yang, S., et al. (2025). Why do multi-agent LLM
systems fail? \emph{arXiv preprint arXiv:2503.13657}.
\url{https://doi.org/10.48550/arXiv.2503.13657}

Chari, A., Tiwari, A., Lian, R., Reddy, S., \& Zhou, B. (2025).
Pheromone-based learning of optimal reasoning paths. \emph{arXiv
preprint arXiv:2501.19278}.
\url{https://doi.org/10.48550/arXiv.2501.19278}

Chen, Z., Zhao, Z., Zhou, A., Huang, Y., Wang, D., Poon, J., \& Li, J.
(2024). AgentPoison: Red-teaming LLM agents via poisoning memory or
knowledge bases. In \emph{Advances in Neural Information Processing
Systems 37 (NeurIPS 2024)}. \url{https://arxiv.org/abs/2407.12784}

CodeRosetta. (2024). Pushing the boundaries of unsupervised code
translation for parallel programming. \emph{Advances in Neural
Information Processing Systems 37 (NeurIPS 2024)}.

Cognitive RPA Framework. (2025). Cognitive RPA: A framework for
hybridizing artificial intelligence with robotic process automation.
\emph{European Journal of Computer Science and Information Technology,
13}(19), 64-78.

Creswell, J. W., \& Creswell, J. D. (2018). \emph{Research design:
Qualitative, quantitative, and mixed methods approaches} (5th ed.). SAGE
Publications.

Cursor. (2025). Scaling long-running autonomous coding. \emph{Cursor
Blog}. \url{https://cursor.com/blog/scaling-agents}

Desai, A., et al. (2025). What is reproducibility in artificial
intelligence and machine learning research? \emph{arXiv preprint
arXiv:2407.10239}. \url{https://doi.org/10.48550/arXiv.2407.10239}

Dietvorst, B. J., Simmons, J. P., \& Massey, C. (2015). Algorithm
aversion: People erroneously avoid algorithms after seeing them err.
\emph{Journal of Experimental Psychology: General, 144}(1), 114-126.
\url{https://doi.org/10.1037/xge0000033}

Dietvorst, B. J., Simmons, J. P., \& Massey, C. (2018). Overcoming
algorithm aversion: People will use algorithms if they can (even
slightly) modify them. \emph{Management Science, 64}(3), 1155-1170.
\url{https://doi.org/10.1287/mnsc.2016.2643}

Dig, D., \& Johnson, R. (2005). The role of refactorings in API
evolution. In \emph{IEEE International Conference on Software
Maintenance (ICSM 2005)} (pp. 2-11).
\url{https://doi.org/10.1109/ICSM.2005.26}

Diggs, C., Doyle, M., Madan, A., \& Scott, E. O. (2025). Leveraging LLMs
for legacy code modernization: Evaluation of LLM-generated
documentation. In \emph{2025 IEEE/ACM International Workshop on Large
Language Models for Code (LLM4Code@ICSE 2025)}.
\url{https://doi.org/10.1109/LLM4Code66737.2025.00027}

Dooley, K. (1997). A complex adaptive systems model of organization
change. \emph{Nonlinear Dynamics, Psychology, and Life Sciences, 1}(1),
69-97. \url{https://doi.org/10.1023/A:1022375910940}

Dumas, M., Leopold, H., Mendling, J., Rinderle-Ma, S., Smedt, J. D., \&
Zimoch, M. (2026). Agentic business process management systems.
\emph{arXiv preprint arXiv:2601.18833}.
\url{https://doi.org/10.48550/arXiv.2601.18833}

Fink, M. (2025). Human oversight under Article 14 of the EU AI Act.
\emph{SSRN}. \url{https://ssrn.com/abstract=5147196}

Feldman, M. S., \& Pentland, B. T. (2003). Reconceptualizing
organizational routines as a source of flexibility and change.
\emph{Administrative Science Quarterly, 48}(1), 94-118.
\url{https://doi.org/10.2307/3556620}

Floridi, L. (2016). Faultless responsibility: On the nature and
allocation of moral responsibility for distributed moral actions.
\emph{Philosophical Transactions of the Royal Society A, 374}(2083),
20160112. \url{https://doi.org/10.1098/rsta.2016.0112}

Gao, M., Li, Y., Liu, B., et al. (2025). Single-agent or multi-agent
systems? Why not both? \emph{arXiv preprint arXiv:2505.18286}.
\url{https://doi.org/10.48550/arXiv.2505.18286}

García-Ruiz, P., \& Rocchi, M. (2025). Can work remain meaningful under
algorithmic management? A MacIntyrean analysis. \emph{Business Ethics
Quarterly, 35}(1), 1-28. \url{https://doi.org/10.1017/beq.2024.36}

Gelernter, D. (1985). Generative communication in Linda. \emph{ACM
Transactions on Programming Languages and Systems, 7}(1), 80-112.
\url{https://doi.org/10.1145/2363.2433}

Geng, L., \& Chang, E. Y. (2025). REALM-Bench: A benchmark for
evaluating multi-agent systems on real-world, dynamic planning and
scheduling tasks. \emph{arXiv preprint arXiv:2502.18836}.
\url{https://doi.org/10.48550/arXiv.2502.18836}

Ghosh Paul, D., Zhu, H., \& Bayley, I. (2024). Benchmarks and metrics
for evaluations of code generation: A critical review. \emph{arXiv
preprint arXiv:2406.12655}.
\url{https://doi.org/10.48550/arXiv.2406.12655}

Grassé, P.-P. (1959). La reconstruction du nid et les coordinations
interindividuelles chez \emph{Bellicositermes natalensis} et
\emph{Cubitermes} sp. la théorie de la stigmergie. \emph{Insectes
Sociaux, 6}(1), 41-80. \url{https://doi.org/10.1007/BF02223791}

Gregor, S., \& Hevner, A. R. (2013). Positioning and presenting design
science research for maximum impact. \emph{MIS Quarterly, 37}(2),
337-355. \url{https://doi.org/10.25300/MISQ/2013/37.2.01}

Gregor, S., Chandra Kruse, L., \& Seidel, S. (2020). The anatomy of a
design principle. \emph{Journal of the Association for Information
Systems, 21}(6), 1622-1652.
\url{https://doi.org/10.17705/1jais.00649}

Grisold, T., Berente, N., \& Seidel, S. (2025). Guardrails for human-AI
ecologies: Norm-based coordination and design for predictability.
\emph{MIS Quarterly, 49}(4), 1239-1266.
\url{https://doi.org/10.25300/MISQ/2025/18058}

Heylighen, F. (2016a). Stigmergy as a universal coordination mechanism
I: Definition and components. \emph{Cognitive Systems Research, 38},
4-13. \url{https://doi.org/10.1016/j.cogsys.2015.12.002}

Heylighen, F. (2016b). Stigmergy as a universal coordination mechanism
II: Varieties and evolution. \emph{Cognitive Systems Research, 38},
50-59. \url{https://doi.org/10.1016/j.cogsys.2015.12.007}

Holland, J. H. (1995). \emph{Hidden order: How adaptation builds
complexity}. Addison-Wesley.

Holmström, J., et al. (2023). \textquotesingle Human
oversight\textquotesingle{} in the EU artificial intelligence act: What,
when and by whom? \emph{Law, Innovation and Technology} (Taylor \&
Francis). \url{https://doi.org/10.1080/17579961.2023.2293129}

Hevner, A. R., March, S. T., Park, J., \& Ram, S. (2004). Design
science in information systems research. \emph{MIS Quarterly, 28}(1),
75-105. \url{https://doi.org/10.2307/25148625}

Hong, S., Zhuge, M., Chen, J., et al. (2024). MetaGPT: Meta programming
for a multi-agent collaborative framework. \emph{arXiv preprint
arXiv:2308.00352}. \url{https://doi.org/10.48550/arXiv.2308.00352}

IBM. (2023). \emph{IBM unveils watsonx generative AI capabilities to
accelerate mainframe application modernization}.
\url{https://newsroom.ibm.com/2023-08-22-IBM-Unveils-watsonx-Generative-AI-Capabilities-to-Accelerate-Mainframe-Application-Modernization}

Jarrahi, M. H., \& Ritala, P. (2025, July 23). Rethinking AI agents: A
principal-agent perspective. \emph{California Management Review}
(Insight).
\url{https://cmr.berkeley.edu/2025/07/rethinking-ai-agents-a-principal-agent-perspective/}

Jimenez, C. E., Yang, J., Wettig, A., et al. (2024). SWE-bench: Can
language models resolve real-world GitHub issues? \emph{arXiv preprint
arXiv:2310.06770}. \url{https://doi.org/10.48550/arXiv.2310.06770}

Kapoor, S., Stroebl, B., Siegel, Z. S., Nadgir, N., \& Narayanan, A.
(2024). AI agents that matter. \emph{arXiv preprint arXiv:2407.01502}.
\url{https://doi.org/10.48550/arXiv.2407.01502}

Kauffman, S. A. (1993). \emph{The origins of order: Self-organization
and selection in evolution}. Oxford University Press.

Kim, H.-W., \& Kankanhalli, A. (2009). Investigating user resistance to
information systems implementation: A status quo bias perspective.
\emph{MIS Quarterly, 33}(3), 567-582.
\url{https://doi.org/10.2307/20650309}

Lamothe, M., Guéhéneuc, Y.-G., \& Shang, W. (2021). A systematic review
of API evolution literature. \emph{ACM Computing Surveys, 54}(8),
171:1-171:36. \url{https://doi.org/10.1145/3450268}

Lee, D., \& Tiwari, M. (2024). Prompt injection: LLM-to-LLM prompt
injection within multi-agent systems. \emph{arXiv preprint
arXiv:2410.07283}. \url{https://arxiv.org/abs/2410.07283}

Lee, J. D., \& See, K. A. (2004). Trust in automation: Designing for
appropriate reliance. \emph{Human Factors, 46}(1), 50-80.
\url{https://doi.org/10.1518/hfes.46.1.50_30392}

Liu, X., Yu, H., Zhang, H., et al. (2024). AgentBench: Evaluating LLMs
as agents. \emph{arXiv preprint arXiv:2308.03688}.
\url{https://doi.org/10.48550/arXiv.2308.03688}

Mamei, M., \& Zambonelli, F. (2009). Programming pervasive and mobile
computing applications: The TOTA approach. \emph{ACM Transactions on
Software Engineering and Methodology, 18}(4), 1-56.
\url{https://doi.org/10.1145/1538942.1538945}

Malone, T. W., \& Crowston, K. (1994). The interdisciplinary study of
coordination. \emph{ACM Computing Surveys, 26}(1), 87-119.
\url{https://doi.org/10.1145/174666.174668}

Markus, M. L. (1983). Power, politics, and MIS implementation.
\emph{Communications of the ACM, 26}(6), 430-444.
\url{https://doi.org/10.1145/358141.358148}

Marsh, L., \& Onof, C. (2008). Stigmergic epistemology, stigmergic
cognition. \emph{Cognitive Systems Research, 9}(1-2), 136-149.
\url{https://doi.org/10.1016/j.cogsys.2007.06.009}

Miles, M. B., Huberman, A. M., \& Saldaña, J. (2020). \emph{Qualitative
data analysis: A methods sourcebook} (4th ed.). SAGE Publications.

Mukherjee, A., \& Chang, H. (2025). Operational agency: A framework for
tracing intent and liability in multi-agent artificial intelligence
systems. \emph{SSRN Electronic Journal}.
\url{https://doi.org/10.2139/ssrn.5136065}

Parunak, H. V. D. (1997). "Go to the ant": Engineering principles from
natural agent systems. \emph{Annals of Operations Research, 75}, 69-101.
\url{https://doi.org/10.1023/A:1018980001403}

OpenAI. (2026, February 5). \emph{Introducing GPT‑5.3‑Codex}.
\url{https://openai.com/index/introducing-gpt-5-3-codex/}

OpenAI. (2022, November 30). \emph{Introducing ChatGPT}.
\url{https://openai.com/index/chatgpt/}

OpenAI. (2026b). \emph{Codex CLI Documentation}.
\url{https://developers.openai.com/codex/cli}

Peng, S., Kalliamvakou, E., Cihon, P., \& Demirer, M. (2023). The impact
of AI on developer productivity: Evidence from GitHub Copilot.
\emph{arXiv preprint arXiv:2302.06590}.
\url{https://doi.org/10.48550/arXiv.2302.06590}

Peffers, K., Tuunanen, T., Rothenberger, M. A., \& Chatterjee, S.
(2007). A design science research methodology for information systems
research. \emph{Journal of Management Information Systems, 24}(3),
45-77. \url{https://doi.org/10.2753/MIS0742-1222240302}

Rahman, M. A. U., Schranz, M., \& Hayat, S. (2025). LLM-powered swarms:
A new frontier or a conceptual stretch? \emph{arXiv preprint
arXiv:2506.14496}. \url{https://doi.org/10.48550/arXiv.2506.14496}

Ricci, A., Omicini, A., Viroli, M., Gardelli, L., \& Oliva, E. (2007).
Cognitive stigmergy: Towards a framework based on agents and artifacts.
In D. Weyns, H. V. D. Parunak, \& F. Michel (Eds.), \emph{Environments
for Multi-Agent Systems III} (LNCS Vol. 4389, pp. 124-140). Springer.
\url{https://doi.org/10.1007/978-3-540-71103-2_7}

Rodriguez, R. R., Jr. (2026). Emergent coordination in multi-agent
systems via pressure fields and temporal decay. \emph{arXiv preprint
arXiv:2601.08129}. \url{https://doi.org/10.48550/arXiv.2601.08129}

Rozière, B., Lachaux, M.-A., Chanussot, L., \& Lample, G. (2020).
Unsupervised translation of programming languages. \emph{Advances in
Neural Information Processing Systems 33 (NeurIPS 2020)}, 20601-20611.
\url{https://doi.org/10.48550/arXiv.2009.08732}

Santoni de Sio, F., \& van den Hoven, J. (2018). Meaningful human
control over autonomous systems: A philosophical account.
\emph{Frontiers in Robotics and AI, 5}, Article 15.
\url{https://doi.org/10.3389/frobt.2018.00015}

Serugendo, G. D. M., Gleizes, M. P., \& Karageorgos, A. (2005).
Self-organization in multi-agent systems. \emph{The Knowledge
Engineering Review, 20}(2), 165-189.
\url{https://doi.org/10.1017/S0269888905000494}

Shrestha, Y. R., Ben-Menahem, S. M., \& von Krogh, G. (2019).
Organizational decision-making structures in the age of artificial
intelligence. \emph{California Management Review, 61}(4), 66-83.
\url{https://doi.org/10.1177/0008125619862257}

Sneed, H. M. (2010). Migrating from COBOL to Java. In \emph{IEEE
International Conference on Software Maintenance (ICSM 2010)} (pp.
237-246). \url{https://doi.org/10.1109/ICSM.2010.5609741}

Strong, D. M., Volkoff, O., Johnson, S. A., Pelletier, L. R., Tulu, B.,
Bar-On, I., Trudel, J., \& Garber, L. (2014). A theory of
organization-EHR affordance actualization. \emph{Journal of the
Association for Information Systems, 15}(2), 53-85.
\url{https://doi.org/10.17705/1jais.00353}

Teece, D. J. (2007). Explicating dynamic capabilities: The nature and
microfoundations of (sustainable) enterprise performance.
\emph{Strategic Management Journal, 28}(13), 1319-1350.
\url{https://doi.org/10.1002/smj.640}

Theraulaz, G., \& Bonabeau, E. (1999). A brief history of stigmergy.
\emph{Artificial Life, 5}(2), 97-116.
\url{https://doi.org/10.1162/106454699568700}

Turel, O., \& Kalhan, S. (2023). Prejudiced against the machine?
Implicit associations and the transience of algorithm aversion.
\emph{MIS Quarterly, 47}(4), 1369-1394.
\url{https://doi.org/10.25300/MISQ/2023/16860}

Uhl-Bien, M., Marion, R., \& McKelvey, B. (2007). Complexity leadership
theory: Shifting leadership from the industrial age to the knowledge
era. \emph{The Leadership Quarterly, 18}(4), 298-318.
\url{https://doi.org/10.1016/j.leaqua.2007.04.002}

Uluhan, E., \& Aydin, M. N. (2014). Complex adaptive systems theory in
the context of business process management. In \emph{S-BPM ONE 2014:
Scientific Research} (pp. 11-23). Springer.
\url{https://doi.org/10.1007/978-3-319-06191-7_2}

Vidgof, M., Bachhofner, S., \& Mendling, J. (2023). Large language
models for business process management: Opportunities and challenges.
\emph{arXiv preprint arXiv:2304.04309}.
\url{https://doi.org/10.48550/arXiv.2304.04309}

Vu, H., Klievtsova, N., Leopold, H., Rinderle-Ma, S., \& Kampik, T.
(2026). Agentic business process management: Practitioner perspectives
on agent governance in business processes. In \emph{BPM 2025 Responsible
BPM Forum. Lecture Notes in Business Information Processing} (vol. 565).
Springer. \url{https://doi.org/10.48550/arXiv.2504.03693}

Venable, J., Pries-Heje, J., \& Baskerville, R. (2016). FEDS: A
framework for evaluation in design science research. \emph{European
Journal of Information Systems, 25}(1), 77-89.
\url{https://doi.org/10.1057/ejis.2014.36}

Wang, X., Li, B., Song, Y., et al. (2025). OpenHands: An open platform
for AI software developers as generalist agents. \emph{arXiv preprint
arXiv:2407.16741}. \url{https://doi.org/10.48550/arXiv.2407.16741}

Weyns, D., Omicini, A., \& Odell, J. (2007). Environment as a first
class abstraction in multiagent systems. \emph{Autonomous Agents and
Multi-Agent Systems, 14}(1), 5-30.
\url{https://doi.org/10.1007/s10458-006-0012-0}

Winfield, A. F. T., et al. (2025). On the ethical governance of swarm
robotic systems in the real world. \emph{Philosophical Transactions of
the Royal Society A, 383}(2289), 20240142.
\url{https://doi.org/10.1098/rsta.2024.0142}

Xia, C. S., Deng, Y., Dunn, S., \& Zhang, L. (2024). Agentless:
Demystifying LLM-based software engineering agents. \emph{arXiv preprint
arXiv:2407.01489}. \url{https://doi.org/10.48550/arXiv.2407.01489}

Xie, J., Zhang, K., Chen, J., Zhu, T., Lou, R., Tian, Y., Xiao, Y.,
\& Su, Y. (2024). TravelPlanner: A benchmark for real-world planning
with language agents. In \emph{Proceedings of the 41st International
Conference on Machine Learning (ICML 2024)}.
\url{https://doi.org/10.48550/arXiv.2402.01622}

Yan, W. (2025, June 12). Don\textquotesingle t build multi-agents.
\emph{Cognition AI Blog}. \url{https://www.cognition-labs.com/blog}

Yang, J., Jimenez, C. E., Yao, S., et al. (2024). SWE-agent:
Agent-computer interfaces enable automated software engineering.
\emph{arXiv preprint arXiv:2405.15793}.
\url{https://doi.org/10.48550/arXiv.2405.15793}

Zhang, Y., Lin, C., Tang, S., Chen, H., Zhou, S., Ma, Y., \& Tresp, V.
(2025). SwarmAgentic: Towards fully automated agentic system generation
via swarm intelligence. \emph{arXiv preprint arXiv:2506.15672}.
\url{https://doi.org/10.48550/arXiv.2506.15672}

Zhu, K., Du, H., Hong, Z., et al. (2025). MultiAgentBench: Evaluating
the collaboration and competition of LLM agents. \emph{arXiv preprint
arXiv:2503.01935}. \url{https://doi.org/10.48550/arXiv.2503.01935}

Ziftci, C., Zheng, J., Choudhary, S., Rossbach, C., Chen, W., Li, W.,
Wang, J., \& Zhang, L. (2025). Migrating code at scale with LLMs at
Google. In \emph{Proceedings of the ACM International Conference on the
Foundations of Software Engineering (FSE 2025)}.

\endgroup
